% !PRINTED 616
\weberpage{616}
%<<CONT
a prime, and which itself, if it is not the unit group, is a transitive linear group.

But this is shown immediately by the representation of the substitutions of the group $L$ by the formula (5).

If in it $e=1$, then $L$ is the cyclic group of degree $n$ with the normal divisor 1. But if $e>1$, let $p$ be a prime dividing $e$, and $e=pe'$. The linear group $L'$ of degree $ne'$, which consists of the substitutions:

\[
\lambda' = (z, a_0^{ph} z + b), \quad
\begin{aligned}
h &= 0, 1, 2 \ldots e'-1\\
b &= 0, 1, 2 \ldots n-1,
\end{aligned}
\]

is certainly a divisor of $L$ of index $p$. But it is also normal, as one recognises without difficulty from the rule of composition (4).

The concepts of the transitive linear and of the metacyclic groups thus coincide completely for the equations of prime degree, and we can therefore in what follows also use the two expressions synonymously.

If we wish to pass from a group $P$ to a conjugate group $P' = \pi^{-1}P\pi$, this transition is effected by carrying out the permutation $\pi$ in the cycles of all the permutations of $P$. $P'$ will therefore agree with $P$ up to the designation of the roots, and is therefore, if $P$ is linear, also to be called linear, even though a change in the numbering of the roots is necessary in order to represent it by linear substitutions.

To simplify the application, let it further be observed that one can derive the \gesperrt{full linear group} by repetition and composition of the two substitutions

\weq{9}{s = (z, z+1), \quad t = (z, gz),}

and every divisor $L$ of the full linear group likewise from

\weq{10}{s = (z, z+1), \quad t^{\alpha_0} = (z, a_0 z)}

One therefore calls the two substitutions (9) or (10) the \gesperrt{generating substitutions of these groups}.

If we set $a_0 = g^2$, we obtain a frequently occurring linear group $(z, az+b)$, in which $a$ runs only through the quadratic residues of $n$, and which has been called by \gesperrt{Kronecker} the \gesperrt{half-metacyclic group}.

This representation by the generating substitutions admits a simple inference concerning the relation of the meta-
%CONT>>
